\section{Introduction}

Universities devote substantial resources to moving research into patents, licenses, and new firms. The rules that govern
this activity are written into institutional intellectual-property (IP) policies, which specify who owns inventions, what
must be disclosed, how income is shared, how conflicts of interest are handled, and what inventors can expect from the
technology-transfer office (TTO). These documents change. Universities revise them in response to legal developments,
governance priorities, organizational experience, and changes in how they manage commercialization. A recurring question
for universities and state systems is therefore straightforward: when an institution rewrites its IP policy, what changes
in the flow of disclosures, patenting, and licensing that follows?

The literature gives little direct guidance. The best-identified evidence concerns large changes in the ownership regime,
most prominently Norway's abolition of the ``professor's privilege,'' which moved ownership of university inventions from
researchers to institutions and was followed by a large decline in university-researcher startups and patenting
\citep{HvideJones2018}. Cross-institution studies relate specific contractual terms, especially inventor royalty shares, to
commercialization \citep{LachSchankerman2008,ArqueCastellsEtAl2016}, although these relationships weaken when policy data
are corrected \citep{OuelletteTutt2020}. Work on technology-transfer offices instead emphasizes staffing, routines, and
relationships \citep{SiegelWaldmanLink2003,MarkmanEtAl2005,DebackereVeugelers2005}. What is missing is systematic evidence
on the ordinary revisions that universities actually undertake, which rarely change ownership and often alter several
provisions and the language surrounding them at once.

This paper studies those revisions. It uses the Policy Communication Stance Index (PCSI), developed and validated in a
companion paper \citep{SharmaPCSI2026}, to characterize whether a revision moves policy language toward a more supportive
or a more restrictive stance toward inventors. Rather than relying primarily on annual policy-in-force scores, most of
which are carried forward from older documents, we focus on moments at which one observed policy document replaces
another. We reconstruct the observed document sequence at each institution, identify changes in PCSI larger than 0.03,
and use adoption dates and revision histories printed in the documents to date the changes.

A central contribution of the revised analysis is a complete and auditable sample-construction exercise. At the baseline
threshold, the linked policy histories contain 127 revisions at 78 universities. We first apply outcome-support and
non-overlap rules that do not use the sign or magnitude of the technology-transfer estimates; 62 revisions remain
mechanically eligible for an event study. We then manually review the predecessor and successor documents for all 62
eligible events. Thirty-four revisions meet the paper's documentary criteria: the pair is comparable or partially
comparable and the revision date is sufficiently documented. The preferred sample therefore contains 34 hand-reviewed
events at 32 institutions, eight upward and 26 downward. A stricter sensitivity sample requiring at least one observed
calendar year after the event contains 29 revisions, seven upward and 22 downward.

The design is a stacked event study \citep{CengizEtAl2019,BakerLarckerWang2022}. Each revising university is compared with
contemporaneous universities that do not revise anywhere in the same event window, with comparisons further saturated by
research-size tercile and public/private status. The headline estimand is the post-revision difference between upward and
downward revisions relative to their clean controls. This comparison nets out changes common to policy revision as such,
but it does not make revision direction random. Universities choose whether and how to revise, so the analysis remains
observational.

The central finding concerns licensing. In the preferred 34-event sample, the average post-revision difference in log
licenses and options between upward and downward revisions is 0.495 (clustered SE 0.140). A 100,000-draw permutation test
that preserves the eight upward labels gives $p=0.032$, and a joint test of the three pre-revision direction coefficients
gives $p=0.32$. The estimate reflects movement in both directions: upward revisions are followed by a 0.279 log-point
increase relative to controls, while downward revisions are followed by a 0.216 log-point decline. The stricter 29-event
sample produces an almost identical licensing gap of 0.499 (SE 0.133; permutation $p=0.026$; pre-trend $p=0.52$).

Other outcomes are less precise. The filing margin and new patent applications move in the same direction as licensing,
but their permutation tests do not reject zero. Invention disclosures do not rise after upward revisions, and patents
issued show no systematic differential response. This pattern locates the strongest post-revision change downstream of
faculty entry into the formal technology-transfer process. At the same time, adjustment for the five outcome tests raises
the licensing Benjamini--Hochberg $q$-value to 0.159. We therefore treat the licensing result as a robust, outcome-specific
empirical pattern rather than as a general rejection across the technology-transfer pipeline.

Several checks sharpen that interpretation. Requiring post-event outcome support reduces the sample from 34 to 29 but
leaves the licensing estimate essentially unchanged. Returning to the original September-25 25-event specification
produces a larger gap of 0.621 in our current replication, while a much broader 59-event specification that admits
rule-coded revisions produces a smaller but still positive estimate. Dropping each event from the 34-event sample in turn
leaves the licensing gap between 0.436 and 0.559. Upward and downward revisions also have nearly identical pre-revision
licensing trends; the clearest difference remains their starting PCSI, with upward revisions beginning from substantially
less supportive prior documents.

The paper contributes evidence on a setting between cross-sectional comparisons of individual policy terms and large
national ownership reforms: ordinary within-university policy revisions. It also contributes a transparent documentary
workflow for turning archival policy histories into dated empirical events. The limitations are consequential. There are
only eight upward revisions in the preferred sample; direction is chosen rather than assigned; the exchangeability
assumption behind permutation inference is substantive rather than design-based; and some policy histories remain
incomplete. The evidence therefore documents a stable post-revision pattern, not a causal effect of particular words or
clauses.

Section~\ref{sec:lit} develops the argument from the literature. Section~\ref{sec:data} describes the policy and AUTM data
and the revision-screening procedure. Section~\ref{sec:strategy} presents the stacked design and inference.
Section~\ref{sec:results} reports the results and robustness. Section~\ref{sec:discussion} discusses interpretation and
limitations, and Section~\ref{sec:conclusion} concludes.
